% GENERATED FILE. Edit canonical lessons, then run npm run generate:books.
% canonical-source-hash: fnv1a64:e62f60afee2fc03e
% canonical-lessons: TA-C387-oppantam, TA-C387-utanpatikkai, TA-C387-vakkuruti, TA-C387-urutimoli, TA-C387-niraiverru

\chapter{Promises and Agreements}
\label{ch:ta-promises-and-agreements}
\hlchaptermodality{387}

\begin{chapteropening}
\textbf{By the end of this chapter:} \emph{I can talk about contracts, promises and pledges.}

Complete the last lesson of chapter 387: I can talk about contracts, promises and pledges.
\end{chapteropening}

\section[\texorpdfstring{\ta{ஒப்பந்தம்} (oppantam)}{ஒப்பந்தம் (oppantam)}]{\ta{ஒப்பந்தம்} (oppantam) — a contract; an agreement}

\label{lesson:TA-C387-oppantam}



\begin{quote}
Before the new one: say the Tamil for to take forward, then the Tamil for an aim.
\end{quote}

\subsection*{You'll want to know: \ta{ஒப்பந்தம்}}
\textbf{\ta{ஒப்பந்தம்}} — \emph{oppantam} — \textquotedblleft{}a contract; an agreement\textquotedblright{}.

\begin{quote}
\ta{இரண்டு} \ta{நிறுவனங்களும்} \ta{ஒப்பந்தத்தில்} \ta{கையெழுத்திட்டன}.
\end{quote}

\begin{quote}
\emph{iraṇṭu niṟuvaṉaṅkaḷum oppantattil kaiyeḻuttiṭṭaṉa.}
\end{quote}

\begin{quote}
The two companies signed the contract.
\end{quote}

\ta{ஒப்பந்தம்} \ta{செய்}, to make an agreement.

\begin{cousinweb}
\textbf{\ta{ஒப்பந்தம்}} (\emph{oppantam}) is native Dravidian: DEDR 924 lists Tamil \emph{o}, \textquotedblleft{}to be suited, acceptable, in concord\textquotedblright{}, with \emph{oppu}, \textquotedblleft{}agreement\textquotedblright{}, and \emph{oppantam}, \textquotedblleft{}agreement\textquotedblright{}.
\end{cousinweb}

\begin{grammarlens}[title={What you've built}]
One more word for saying what happens if.
\end{grammarlens}

\subsection*{Your turn}
\begin{itemize}
\raggedright
  \item \emph{Say it:} \emph{oppantam}
  \item \emph{Say it:} \emph{oppantam}, once more
  \item \emph{Say it:} \emph{iraṇṭu niṟuvaṉaṅkaḷum oppantattil kaiyeḻuttiṭṭaṉa.} — The two companies signed the contract.
  \item \emph{Recall:} say the Tamil for to take forward, then the Tamil for an aim, then say \emph{oppantam} again
\end{itemize}

\begin{tcolorbox}[breakable,colback=teal!4,colframe=teal!35!black,title={Before you move on}]
What does \ta{ஒப்பந்தம்} mean? (\textquotedblleft{}A contract; an agreement\textquotedblright{}.) Say it once more.
\end{tcolorbox}
\section[\texorpdfstring{\ta{உடன்படிக்கை} (uṭaṉpaṭikkai)}{உடன்படிக்கை (uṭaṉpaṭikkai)}]{\ta{உடன்படிக்கை} (uṭaṉpaṭikkai) — a treaty; a pact}

\label{lesson:TA-C387-utanpatikkai}



\begin{quote}
Before the new one: say the Tamil for an aim, then the Tamil for a contract.
\end{quote}

\subsection*{You'll want to know: \ta{உடன்படிக்கை}}
\textbf{\ta{உடன்படிக்கை}} — \emph{uṭaṉpaṭikkai} — \textquotedblleft{}a treaty; a pact\textquotedblright{}.

\begin{quote}
\ta{இரண்டு} \ta{நாடுகளும்} \ta{அமைதி} \ta{உடன்படிக்கை} \ta{செய்தன}.
\end{quote}

\begin{quote}
\emph{iraṇṭu nāṭukaḷum amaiti uṭaṉpaṭikkai ceytaṉa.}
\end{quote}

\begin{quote}
The two countries made a peace treaty.
\end{quote}

Between countries; \ta{ஒப்பந்தம்} is the general word.

\begin{cousinweb}
\textbf{\ta{உடன்படிக்கை}} (\emph{uṭaṉpaṭikkai}) is native, built on \emph{uṭaṉ}, \textquotedblleft{}together\textquotedblright{}, which DEDR 945 lists with \emph{uṭaṉ-paṭu}, \textquotedblleft{}to agree, yield\textquotedblright{}: a written agreeing.
\end{cousinweb}

\begin{grammarlens}[title={What you've built}]
One more word for saying what happens if.
\end{grammarlens}

\subsection*{Your turn}
\begin{itemize}
\raggedright
  \item \emph{Say it:} \emph{uṭaṉpaṭikkai}
  \item \emph{Say it:} \emph{uṭaṉpaṭikkai}, once more
  \item \emph{Say it:} \emph{iraṇṭu nāṭukaḷum amaiti uṭaṉpaṭikkai ceytaṉa.} — The two countries made a peace treaty.
  \item \emph{Recall:} say the Tamil for an aim, then the Tamil for a contract, then say \emph{uṭaṉpaṭikkai} again
\end{itemize}

\begin{tcolorbox}[breakable,colback=teal!4,colframe=teal!35!black,title={Before you move on}]
What does \ta{உடன்படிக்கை} mean? (\textquotedblleft{}A treaty; a pact\textquotedblright{}.) Say it once more.
\end{tcolorbox}
\section[\texorpdfstring{\ta{வாக்குறுதி} (vākkuṟuti)}{வாக்குறுதி (vākkuṟuti)}]{\ta{வாக்குறுதி} (vākkuṟuti) — a promise}

\label{lesson:TA-C387-vakkuruti}



\begin{quote}
Before the new one: say the Tamil for a contract, then the Tamil for a treaty.
\end{quote}

\subsection*{You'll want to know: \ta{வாக்குறுதி}}
\textbf{\ta{வாக்குறுதி}} — \emph{vākkuṟuti} — \textquotedblleft{}a promise\textquotedblright{}.

\begin{quote}
\ta{தேர்தலுக்கு} \ta{முன்} \ta{அவர்} \ta{பல} \ta{வாக்குறுதிகள்} \ta{கொடுத்தார்}.
\end{quote}

\begin{quote}
\emph{tērtalukku muṉ avar pala vākkuṟutikaḷ koṭuttār.}
\end{quote}

\begin{quote}
Before the election he made many promises.
\end{quote}

Plural: \ta{வாக்குறுதிகள்}, \emph{vākkuṟutikaḷ}.

\ta{வாக்குறுதி} \ta{கொடு}, to promise; \ta{வாக்குறுதியைக்} \ta{காப்பாற்று}, to keep a promise.

\begin{cousinweb}
\textbf{\ta{வாக்குறுதி}} (\emph{vākkuṟuti}) is built from \emph{vākku}, \textquotedblleft{}word\textquotedblright{}, borrowed from Sanskrit \emph{vāc} (Monier-Williams), and native \emph{uṟuti}, \textquotedblleft{}firmness, assurance\textquotedblright{} (DEDR 721): a firm word.
\end{cousinweb}

\begin{grammarlens}[title={What you've built}]
One more word for saying what happens if.
\end{grammarlens}

\subsection*{Your turn}
\begin{itemize}
\raggedright
  \item \emph{Say it:} \emph{vākkuṟuti}
  \item \emph{Say it:} \emph{vākkuṟuti}, once more
  \item \emph{Say it:} \emph{tērtalukku muṉ avar pala vākkuṟutikaḷ koṭuttār.} — Before the election he made many promises.
  \item \emph{Recall:} say the Tamil for a contract, then the Tamil for a treaty, then say \emph{vākkuṟuti} again
\end{itemize}

\begin{tcolorbox}[breakable,colback=teal!4,colframe=teal!35!black,title={Before you move on}]
What does \ta{வாக்குறுதி} mean? (\textquotedblleft{}A promise\textquotedblright{}.) Say it once more.
\end{tcolorbox}
\section[\texorpdfstring{\ta{உறுதிமொழி} (uṟutimoḻi)}{உறுதிமொழி (uṟutimoḻi)}]{\ta{உறுதிமொழி} (uṟutimoḻi) — a pledge; an oath}

\label{lesson:TA-C387-urutimoli}



\begin{quote}
Before the new one: say the Tamil for a treaty, then the Tamil for a promise.
\end{quote}

\subsection*{You'll want to know: \ta{உறுதிமொழி}}
\textbf{\ta{உறுதிமொழி}} — \emph{uṟutimoḻi} — \textquotedblleft{}a pledge; an oath\textquotedblright{}.

\begin{quote}
\ta{மாணவர்கள்} \ta{ஊரைச்} \ta{சுத்தமாக} \ta{வைப்போம்} \ta{என்று} \ta{உறுதிமொழி} \ta{எடுத்தார்கள்}.
\end{quote}

\begin{quote}
\emph{māṇavarkaḷ ūraic cuttamāka vaippōm eṉṟu uṟutimoḻi eṭuttārkaḷ.}
\end{quote}

\begin{quote}
The students took a pledge to keep the town clean.
\end{quote}

\ta{உறுதிமொழி} \ta{எடு}, to take a pledge.

\begin{cousinweb}
\textbf{\ta{உறுதிமொழி}} (\emph{uṟutimoḻi}) is native: \emph{uṟuti}, \textquotedblleft{}firmness, certainty\textquotedblright{} (DEDR 721), and \emph{moḻi}, \textquotedblleft{}word, saying\textquotedblright{} (DEDR 4989): a firm word, spoken out.
\end{cousinweb}

\begin{grammarlens}[title={What you've built}]
One more word for saying what happens if.
\end{grammarlens}

\subsection*{Your turn}
\begin{itemize}
\raggedright
  \item \emph{Say it:} \emph{uṟutimoḻi}
  \item \emph{Say it:} \emph{uṟutimoḻi}, once more
  \item \emph{Say it:} \emph{māṇavarkaḷ ūraic cuttamāka vaippōm eṉṟu uṟutimoḻi eṭuttārkaḷ.} — The students took a pledge to keep the town clean.
  \item \emph{Recall:} say the Tamil for a treaty, then the Tamil for a promise, then say \emph{uṟutimoḻi} again
\end{itemize}

\begin{tcolorbox}[breakable,colback=teal!4,colframe=teal!35!black,title={Before you move on}]
What does \ta{உறுதிமொழி} mean? (\textquotedblleft{}A pledge; an oath\textquotedblright{}.) Say it once more.
\end{tcolorbox}
\section[\texorpdfstring{\ta{நிறைவேற்று} (niṟaivēṟṟu)}{நிறைவேற்று (niṟaivēṟṟu)}]{\ta{நிறைவேற்று} (niṟaivēṟṟu) — to fulfil; to carry out}

\label{lesson:TA-C387-niraiverru}



\begin{quote}
Before the new one: say the Tamil for a promise, then the Tamil for a pledge.
\end{quote}

\subsection*{You'll want to know: \ta{நிறைவேற்று}}
\textbf{\ta{நிறைவேற்று}} — \emph{niṟaivēṟṟu} — \textquotedblleft{}to fulfil; to carry out\textquotedblright{}.

\begin{quote}
\ta{அவர்} \ta{தன்} \ta{வாக்குறுதியை} \ta{நிறைவேற்றினார்}.
\end{quote}

\begin{quote}
\emph{avar taṉ vākkuṟutiyai niṟaivēṟṟiṉār.}
\end{quote}

\begin{quote}
He fulfilled his promise.
\end{quote}

Past (I …): \ta{நிறைவேற்றினேன்}, \emph{niṟaivēṟṟiṉēṉ}.

Also to pass a law: \ta{சட்டத்தை} \ta{நிறைவேற்று}. Its intransitive partner is \ta{நிறைவேறு}, to come true.

\begin{cousinweb}
\textbf{\ta{நிறைவேற்று}} (\emph{niṟaivēṟṟu}) is native Dravidian: DEDR 3682 lists Tamil \emph{niṟai}, \textquotedblleft{}to become full\textquotedblright{}, and \emph{niṟaivēṟṟu}, \textquotedblleft{}to accomplish\textquotedblright{}: to bring something to fullness.
\end{cousinweb}

\begin{grammarlens}[title={What you've built}]
One more word for saying what happens if.
\end{grammarlens}

\subsection*{Your turn}
\begin{itemize}
\raggedright
  \item \emph{Say it:} \emph{niṟaivēṟṟu}
  \item \emph{Say it:} \emph{niṟaivēṟṟu}, once more
  \item \emph{Say it:} \emph{avar taṉ vākkuṟutiyai niṟaivēṟṟiṉār.} — He fulfilled his promise.
  \item \emph{Recall:} say the Tamil for a promise, then the Tamil for a pledge, then say \emph{niṟaivēṟṟu} again
\end{itemize}

\begin{tcolorbox}[breakable,colback=teal!4,colframe=teal!35!black,title={Before you move on}]
What does \ta{நிறைவேற்று} mean? (\textquotedblleft{}To fulfil; to carry out\textquotedblright{}.) Say it once more.
\end{tcolorbox}
